\chapter{Apprendre par la pratique : \texttt{forge-apprendre}}
\source{skills/forge-apprendre/SKILL.md et lecons/}
Le skill est un \emph{coach} : il explique, mais c'est vous qui écrivez chaque livrable. Les exercices portent sur un cas fictif, \textbf{PayBoutique} (boutique en ligne béninoise qui encaisse des paiements MTN MoMo et Moov Money), et vivent dans \cmd{apprentissage/}, jamais dans \cmd{docs/}.

\section{Les quatre leçons}
\noindent\begin{tabularx}{\linewidth}{lXX}
\toprule
\textbf{Leçon} & \textbf{Exercice} & \textbf{Vérification} \\ \midrule
1 Impact map & objectif mesurable, acteur, tri de trois idées & mots-clés \emph{indicateur}, \emph{cible}, \emph{échéance} \\
2 Langage & glossaire et code avec un terme banni & \cmd{check-glossary-terms.sh} : KO puis OK \\
3 Spec & règle de paiement avec trois cas & \cmd{check-spec-validated.sh} \\
4 ADR & choix d'hébergement chiffré en FCFA & \cmd{check-adr-present.sh} \\
\bottomrule
\end{tabularx}

\section{Suivre sa progression}
\cmd{scripts/forge-lecons.sh} affiche par exemple \cmd{Leçons [██░░] 2/4 · prochaine : 03-spec}, calculé à partir de vos fichiers d'exercice : vous reprenez là où vous vous étiez arrêté, sans fichier de suivi.

\astuce{Les scripts ne vérifient que la présence et le statut, pas la qualité. La qualité se discute en relecture avec le coach. Les concepts sans leçon (bounded context, OpenAPI d'abord, fitness functions, recette) s'étudient à partir du skill correspondant.}

\section{Après chaque leçon}
Refaites l'exercice sur votre projet réel avec le skill de l'étape : \cmd{forge-cadrage}, \cmd{forge-domaine}, \cmd{forge-spec} (ou \cmd{forge-rapide}), \cmd{forge-architecture}.
